\documentclass[11pt,a4paper]{article}
\usepackage[utf8]{inputenc}
\usepackage[T1]{fontenc}
\usepackage[portuguese]{babel}
\usepackage{amsmath}
\usepackage{amsfonts}
\usepackage{amssymb}
\usepackage{amsthm}
\usepackage{mathtools}
\usepackage{graphicx}
\usepackage{booktabs}
\usepackage{hyperref}
\usepackage{cleveref}
\usepackage{geometry}
\usepackage{tikz}
\usepackage{tikz-cd}
\usepackage{listings}
\usepackage{xcolor}

\geometry{margin=2.5cm}

\hypersetup{
    colorlinks=true,
    linkcolor=blue,
    citecolor=red,
    urlcolor=blue
}

\lstset{
    language=Rust,
    basicstyle=\ttfamily\small,
    keywordstyle=\color{blue},
    commentstyle=\color{green!60!black},
    stringstyle=\color{red},
    numbers=left,
    numberstyle=\tiny,
    frame=single,
    breaklines=true
}

\newtheorem{theorem}{Teorema}[section]
\newtheorem{definition}[theorem]{Definição}

\newcommand{\rhoT}{\rho_T}
\newcommand{\BT}{B_T}
\newcommand{\CT}{C_T}
\newcommand{\RT}{R_T}
\newcommand{\DeltaT}{\Delta T}
\newcommand{\RL}{\mathcal{R}_L}

\title{
    \vspace{-1cm}
    \textbf{Arquitetura Completa do MutaCore no Ecossistema Reino Digital}\\
    \large Organismo Cibernético: Lógica, Transporte, Metabolismo e Governança
}
\author{Euzébio Soares \and Reino Digital -- Core Engine}
\date{Outubro de 2026 -- Versão 4.0.26}

\begin{document}

\maketitle

\begin{abstract}
Apresentamos a arquitetura formal completa do MutaCore, o organismo cibernético central do ecossistema Reino Digital. O MutaCore unifica quatro camadas interdependentes: (1) Substrato Lógico com limiares de Krajicek/Gödel, (2) Transporte via Origo Overlay com ASTs alfa-estruturais, (3) Metabolismo Termodinâmico com acoplamento $\DeltaT \leftrightarrow \RL$, e (4) Membrana de Governança com dNFTs e oráculos contra-factuais. Formalizamos a dinâmica evolutiva, provamos propriedades de segurança, e especificamos interfaces de integração.
\end{abstract}

\tableofcontents
\newpage

\section{Visão Geral da Arquitetura}

O MutaCore não é uma arquitetura de microsserviços tradicional. É um \textbf{organismo cibernético completo} onde cada camada mapeia para um sistema biológico análogo:

\begin{center}
\begin{tabular}{@{}lll@{}}
\toprule
\textbf{Camada MutaCore} & \textbf{Análogo Biológico} & \textbf{Primitiva Computacional} \\
\midrule
1. Substrato Lógico & Esqueleto / Rigidez & Limiares de Prova (Krajicek/Gödel) \\
2. Transporte & Sistema Nervoso & Origo Overlay / AST Alfa-Estrutural \\
3. Metabolismo & Metabolismo / Hormônios & Resíduo de Landauer / $\DeltaT$ \\
4. Membrana & Sistema Imune / Pele & dNFT Containers / Q-ESCUDO \\
\bottomrule
\end{tabular}
\end{center}

\subsection{Princípios de Design}

\begin{enumerate}
    \item \textbf{Formal-First:} Toda mutação passa por validação de limiares de prova \textit{antes} da execução.
    \item \textbf{Thermodynamic-Aware:} Custo energético (Landauer) é métrica de primeira classe, não afterthought.
    \item \textbf{Alpha-Structural:} Código é transportado como ASTs canônicas, não texto — imune a renomeação.
    \item \textbf{Distributed-by-Default:} Quantum Local Packing otimiza localidade em hipercubo de nós.
    \item \textbf{Self-Governing:} Identidade e direitos via dNFTs; segurança física via oráculos contra-factuais.
\end{enumerate}

\section{Camada 1: Substrato Lógico e Criptográfico}

\subsection{Espectro de Krajicek}

A base lógica é o espectro $\mathcal{K} = (\rhoT, \BT, \CT, \RT)$ definido em \texttt{krajicek\_thresholds.tex}. 

\begin{theorem}[Segurança Lógica]
Nenhuma mutação com complexidade de prova $> \rhoT$ pode ser aceita pelo sistema.
\end{theorem}

\subsection{Equivariância no Hipercubo}

Estados distribuídos são rastreados no hipercubo $Q_n$. A equivariância garante que todos os nós vejam o mesmo estado lógico sem travas centralizadas.

\begin{equation}
    \text{CustoEquivariância} = \sum_{i=1}^n \lambda_i \cdot \text{dist}(x_i, x_i')
\end{equation}
onde $\lambda_i$ são autovalores do Laplaciano de Hodge (ver \texttt{delta\_t\_topology.tex}).

\subsection{Criptografia SPECTRUM-1024}

Cifrador de bloco 1024-bit com prova de estabilidade bijectiva (ver \texttt{spectrum1024\_tests.py}).

\section{Camada 2: Transporte, Comunicação e Distribuição}

\subsection{Origo Overlay}

Implementado em \texttt{origo-overlay} (Rust). Características:

\begin{itemize}
    \item \textbf{Overlays Declarativos:} Especificação de topologia de rede como dados, não código imperativo.
    \item \textbf{AST Alfa-Estrutural:} Hash estrutural invariante sob $\alpha$-conversão (ver \texttt{hash.rs}).
    \item \textbf{Validação Prévia:} \texttt{OrigoAstValidator} verifica limiares de Krajicek \textit{antes} da execução.
\end{itemize}

\subsection{Quantum Local Packing}

Algoritmo de annealing quântico para alocação ótima de cargas (ver \texttt{quantum\_packing.rs}, \texttt{hodge\_matrices.py}).

\begin{equation}
    \min_{\text{atribuição}} \sum_{i,j} w_{ij} \cdot d_{\text{hypercube}}(node_i, node_j) + \text{balanceamento}
\end{equation}

\section{Camada 3: Metabolismo Termodinâmico (MUTARIC)}

\subsection{Acoplamento Fundamental}

O coração do metabolismo é a equação (ver \texttt{landauer\_coupling.tex}):

\begin{equation}
    \RL = \left( \sum_{i=1}^{n} \DeltaT_i^2 \right) \cdot
    \exp\left( \frac{\rhoT - \BT}{\CT \cdot \RT} \right) \cdot
    \kB T \ln 2
\end{equation}

\subsection{Estados Metabólicos}

\begin{center}
\begin{tabular}{@{}lll@{}}
\toprule
\textbf{Estado} & \textbf{Condição} & \textbf{Comportamento} \\
\midrule
HOMEOSTASIS & $\RL$ baixo, $\DeltaT$ baixo & Operação normal, monitoramento \\
EXPLORATION & $\RL$ baixo, $\DeltaT$ médio & Mutações ousadas, expansão \\
ANXIETY & $\RL$ alto ou $\DeltaT$ alto & Contração, simplificação, safe-mode \\
\bottomrule
\end{tabular}
\end{center}

\subsection{ricemotions: Ponte PIXEL/RIC $\leftrightarrow$ TEOA}

\begin{itemize}
    \item \textbf{PIXEL:} Representação visual/espacial do estado
    \item \textbf{RIC:} Código informacional reduzido (AST comprimida)
    \item \textbf{TEOA:} Teoria Epistemológica de Objetos Abstratos (semântica)
\end{itemize}

O metabolismo processa $\text{PIXEL} \to \text{RIC} \to \text{TEOA}$ e monitora perda de informação (Landauer) em cada etapa.

\section{Camada 4: Membrana, Governança e Identidade}

\subsection{dNFT Containers}

Cada agente/mutação é encapsulado em container governado por dNFT (Dynamic NFT):

\begin{itemize}
    \item \textbf{Identidade Genômica:} Hash da AST + linhagem evolutiva
    \item \textbf{Metadados Imutáveis:} Armazenados on-chain (IPFS + blockchain)
    \item \textbf{Direitos Autorais:} Registro automático de IP ao gerar código
    \item \textbf{Evolução Rastreável:} Grafo de mutações pai $\to$ filho
\end{itemize}

\subsection{Q-ESCUDO: Oráculo Contra-Factual}

Validação de segurança no mundo físico (NR-12 industrial):

\begin{equation}
    \text{Seguro}(a) \iff \forall w \in \mathcal{W}: \text{Counterfactual}(w, a) \models \text{NR-12}
\end{equation}
onde $\mathcal{W}$ é o espaço de mundos possíveis, $a$ a ação proposta.

\section{Dinâmica Evolutiva: Ciclo de Vida}

\begin{center}
\begin{tikzpicture}[node distance=2cm, >=stealth]
\node[draw, rectangle, rounded corners, fill=blue!10] (1) {1. Homeostase Estável};
\node[draw, rectangle, rounded corners, fill=green!10, right of=1] (2) {2. Estímulo Exógeno};
\node[draw, rectangle, rounded corners, fill=yellow!10, right of=2] (3) {3. Mutação Sintática};
\node[draw, rectangle, rounded corners, fill=orange!10, right of=3] (4) {4. Filtro de Seleção};
\node[draw, rectangle, rounded corners, fill=red!10, right of=4] (5) {5. Encapsulamento};

\draw[->, thick] (1) -- (2);
\draw[->, thick] (2) -- (3);
\draw[->, thick] (3) -- (4);
\draw[->, thick] (4) -- (5);
\draw[->, thick, bend right=40] (5) to (1);
\end{tikzpicture}
\end{center}

\subsection{Descrição das Fases}

\begin{enumerate}
    \item \textbf{Homeostase:} Monitora $\RL$ e equivariância. Aguarda estímulo.
    \item \textbf{Estímulo:} Dados via Origo Overlay. $\DeltaT$ mede distância semântica.
    \item \textbf{Mutação:} Gera novas ASTs para resolver $\DeltaT$. Python/Rust.
    \item \textbf{Filtro:} Validador aplica limiares $\mathcal{K}$. Rejeita se $> \rhoT$.
    \item \textbf{Encapsulamento:} Aprova $\to$ assina dNFT $\to$ Quantum Packing $\to$ deploy.
\end{enumerate}

\section{Tabuleiro de Interações}

\begin{center}
\begin{tabular}{@{}lll@{}}
\toprule
\textbf{Atributo} & \textbf{Especificação Técnica} & \textbf{Objetivo} \\
\midrule
Substrato Execução & Rust (core) + Python (flex) & Latência mínima + plasticidade \\
Consenso & Prova Complexidade Krajicek & Evita alucinações lógicas \\
Indicador Metabólico & Landauer Monitoring & Otimiza energia/custo cloud \\
Interface Mercado & dNFTs + Briefings Auto & Monetização descentralizada \\
Segurança Física & Q-ESCUDO NR-12 & Validação contra-factual \\
\bottomrule
\end{tabular}
\end{center}

\section{Implementação: Mapa de Arquivos}

\begin{lstlisting}[caption=Estrutura do Crate Rust origo-overlay]
origo-overlay/
├── Cargo.toml
└── src/
    ├── lib.rs              # Exports públicos
    ├── ast_node.rs         # Definição AstNode, AstNodeKind
    ├── alpha_context.rs    # Canonização alfa, isomorfismo
    ├── ast_validator.rs    # Validação limiares Krajicek
    ├── krajicek_thresholds.rs # Espectro ρ_T, B_T, C_T, R_T
    ├── quantum_packing.rs  # Quantum Local Packing
    ├── hash.rs             # Hash alfa-estrutural, Merkle
    └── error.rs            # Tipos de erro OrigoError
\end{lstlisting}

\begin{lstlisting}[caption=Scripts Python Principais]
src/python/
├── landauer_simulation.py   # Simulação termodinâmica ΔT ↔ ℛₗ
├── hodge_matrices.py        # Decomposição Hodge em Q_n (n>6)
└── spectrum1024_tests.py    # Validação bijectiva SPECTRUM-1024
\end{lstlisting}

\section{Interfaces de Integração}

\subsection{Rust $\leftrightarrow$ Python (PyO3)}

\begin{lstlisting}
#[pyfunction]
fn validate_ast(ast_json: &str, rho_t: usize) -> PyResult<String> {
    let ast: AstNode = serde_json::from_str(ast_json)?;
    let mut validator = OrigoAstValidator::new_simple(rho_t);
    let result = validator.validate(&ast)?;
    Ok(serde_json::to_string(&result)?)
}
\end{lstlisting}

\subsection{WASM Sandbox (Planejado v0.4.0)}

\begin{lstlisting}
// Compilação para WebAssembly
// wasm-pack build --target web
// Execução isolada no navegador/Node.js
\end{lstlisting}

\section{Segurança e Confiabilidade}

\subsection{Modelo de Ameaça}

\begin{itemize}
    \item \textbf{Mutação Maliciosa:} Bloqueada por $\rhoT$ (Camada 1)
    \item \textbf{Corrupção de Transporte:} Detectada por hash alfa-estrutural (Camada 2)
    \item \textbf{Exaustão Energética:} Prevenida por monitoramento $\RL$ (Camada 3)
    \item \textbf{Roubo de IP:} Impossível por dNFT lineage (Camada 4)
    \item \textbf{Falha Física:} Prevenida por Q-ESCUDO (Camada 4)
\end{itemize}

\subsection{Propriedades Garantidas}

\begin{theorem}[Integridade Lógica]
Se $\texttt{Validate}(ast, \mathcal{K}) = \text{Accept}$, então $ast$ é consistente com o estado atual do sistema.
\end{theorem}

\begin{theorem}[Terminabilidade Metabólica]
Para $\gamma > 0$, $\RL(t)$ converge para equilíbrio em tempo finito.
\end{theorem}

\begin{theorem}[Não-Repúdio de Identidade]
dNFT lineage garante rastreabilidade completa de autoria de código.
\end{theorem}

\section{Roadmap de Implementação}

\begin{center}
\begin{tabular}{@{}lll@{}}
\toprule
\textbf{Versão} & \textbf{Foco} & \textbf{Entregáveis} \\
\midrule
v0.3.0 & Core Rust & origo-overlay, validador, thresholds \\
v0.4.0 & WASM Sandbox & Execução isolada, PyO3 bindings \\
v0.5.0 & Hodge n>6 & Matrizes bloco, equivariância testada \\
v0.6.0 & SPECTRUM-1024 CI & Prova bijectiva em pipeline \\
v1.0.0 & Release Estável & MutaCore completo, docs, exemplos \\
\bottomrule
\end{tabular}
\end{center}

\section{Conclusão}

O MutaCore estabelece uma nova classe de arquitetura: \textbf{computação orgânico-formal}. Ao integrar:
\begin{enumerate}
    \item Limites lógicos (Gödel/Krajicek) — \textit{esqueleto}
    \item Transporte imutável (AST alfa) — \textit{sistema nervoso}
    \item Termodinâmica da informação (Landauer) — \textit{metabolismo}
    \item Governança criptográfica (dNFT/Q-ESCUDO) — \textit{membrana}
\end{enumerate}
O Reino Digital deixa de ser repositório estático para se tornar organismo auto-referencial, auto-reparável, em constante \textit{vir-a-ser}.

\bibliographystyle{plain}
\begin{thebibliography}{99}
\bibitem{landauer} R. Landauer, ``Irreversibility and Heat Generation...'', 1961.
\bibitem{krajicek} P. Krajicek, \textit{Bounded Arithmetic...}, Cambridge, 1995.
\bibitem{goedel} K. Gödel, ``Über formal unentscheidbare Sätze...'', 1931.
\bibitem{soares2024} E. Soares, ``MutaCore: Cybernetic Organism...'', Reino Digital, 2024.
\end{thebibliography}

\appendix
\section{Apêndice: Glossário de Símbolos}
\begin{tabular}{@{}cl@{}}
\toprule
$\rhoT, \BT, \CT, \RT$ & Limiares de Krajicek \\
$\DeltaT, \DeltaT_i$ & Gap Topológico Narrativo (total, componentes) \\
$\RL$ & Resíduo de Landauer \\
$\MT, \ME$ & Variedades Organismo / Ambiente \\
$Q_n$ & Hipercubo $n$-dimensional \\
$\laplacian_k$ & Laplaciano de Hodge grau $k$ \\
$\mathcal{E}$ & Score de Equivariância \\
dNFT & Dynamic NFT \\
Q-ESCUDO & Quantum-Enhanced Safety Counterfactual \\
\bottomrule
\end{tabular}

\end{document}